\documentclass[a4paper,10pt]{article}
\usepackage[utf8]{inputenc}
\usepackage[T5]{fontenc}
\usepackage[vietnamese]{babel}
\usepackage[top=1.4cm, bottom=1.4cm, left=1.6cm, right=1.6cm]{geometry}
\usepackage{setspace}
\usepackage{enumitem}
\usepackage{sectsty}
\usepackage{xcolor}
\usepackage{hyperref}
\usepackage{graphicx}
\usepackage{booktabs}
\usepackage{array}
\usepackage{tabularx}
\usepackage{caption}
\setstretch{1}

\sectionfont{\color{blue!60!black}\large}
\subsectionfont{\color{blue!50!black}\normalsize}

\begin{document}

% ── Header thủ công (không dùng \maketitle) ───────────────────────────────
\begin{center}
  {\LARGE\bfseries BÁO CÁO THU THẬP DỮ LIỆU: EDUVISION}\\[0.15cm]
  {\large Data Collection Report --- Giai đoạn 1 \quad\textbf{| Nhóm 1}}
\end{center}
\vspace{-0.2cm}
\noindent\rule{\textwidth}{0.4pt}
\vspace{0.1cm}

\section*{1. Thu thập Dữ liệu --- Data Collection}

Nhóm tiến hành quay video tại phòng học thực tế của Trường Đại học FPT theo hai chiến lược:

\textbf{(a) Wide Shot} --- Camera được đặt tại 3 góc lớp học (góc chéo, góc thẳng trái, góc thẳng phải) ở độ cao 2.2--2.5m, bao phủ toàn bộ không gian, phục vụ bài toán Person Detection \& Tracking trong điều kiện thực tế.

\textbf{(b) Expression Data} --- Các thành viên nhóm ngồi tại 3 vị trí khác nhau (hàng đầu, hàng giữa, hàng cuối) và diễn các trạng thái hành vi. Góc quay được giữ đồng nhất với góc triển khai thực tế để đảm bảo tính nhất quán.

\vspace{0.1cm}
\begin{center}
\begin{tabularx}{0.96\textwidth}{lXrr}
\toprule
\textbf{Loại} & \textbf{Hành vi / Góc quay} & \textbf{Vị trí} & \textbf{Số frames} \\
\midrule
Wide Shot & Góc chéo, Góc thẳng trái, Góc thẳng phải & --- & 681 \\
Expression & focused, drowsy, sleeping, using\_phone, off\_task, & 3 & 654 \\
           & side\_talking, away\_from\_seat, raising\_hand & & \\
\midrule
\multicolumn{3}{l}{\textbf{Tổng cộng}} & \textbf{1.335 frames} \\
\bottomrule
\end{tabularx}
\end{center}

\section*{2. Xử lý Dữ liệu --- Data Processing Pipeline}

Toàn bộ video gốc (.MOV, 30--60 fps) được xử lý qua pipeline tự động gồm các bước:

\begin{itemize}[noitemsep, topsep=1pt]
    \item \textbf{Frame Extraction} --- Trích xuất frame ở 2 fps bằng OpenCV (\texttt{tools/extract\_frames.py}), sinh ra 1.335 frame ảnh .jpg được phân loại tự động theo loại (wide shot / expression), vị trí và hành vi dựa trên cấu trúc thư mục.
    \item \textbf{Auto-Annotation} --- Áp dụng mô hình YOLO11n pretrained để tự động sinh bounding box người trên 681 Wide Shot frames (\texttt{tools/auto\_annotate.py}), tạo ra \textbf{7.806 bounding boxes} (trung bình 11.5 người/frame). Nhãn hành vi cho Expression frames được gán tự động từ tên thư mục.
    \item \textbf{Label Review} --- Nhóm review và hiệu chỉnh annotation thủ công trên nền tảng Roboflow (bổ sung behavior label cho từng bounding box trong Wide Shot).
    \item \textbf{Dataset Split} --- Phân chia train/val/test theo tỉ lệ \textbf{70/20/10} (stratify theo góc camera) bằng \texttt{tools/split\_dataset.py}: 475 train / 135 val / 71 test frames.
    \item \textbf{Lưu trữ} --- Toàn bộ dataset (frames + annotations + splits) được lưu tại HuggingFace: \href{https://huggingface.co/datasets/annghoang/EduVision}{\texttt{annghoang/EduVision}}; source code pipeline tại \href{https://github.com/TuanLe303/EduVision}{\texttt{TuanLe303/EduVision}} (branch: \texttt{data-pipeline}).
\end{itemize}

\section*{3. Dataset Công khai Bổ sung --- External Datasets}

Để tăng độ phong phú và tính tổng quát, nhóm bổ sung 3 dataset công khai:

\begin{itemize}[noitemsep, topsep=1pt]
    \item \textbf{SCB-Dataset} \href{https://github.com/Whiffe/SCB-dataset}{[GitHub]} --- \textbf{7.428 ảnh, 106.830 nhãn, 20 classes} hành vi lớp học (hand-raising, read, write, bow head, talk, using phone...). Sẽ dùng để fine-tune module Behavior Analysis cho các class \texttt{focused}, \texttt{using\_phone}, \texttt{raising\_hand}, \texttt{side\_talking}.

    \item \textbf{Ambient Intelligence Classroom} \href{https://universe.roboflow.com/national-taipei-university-of-technology-jc4sl/ambient-intelligence-classroom}{[Roboflow]} --- \textbf{$\sim$1.350 ảnh} (CC BY 4.0, định dạng YOLOv8) với nhãn khớp hoàn toàn EduVision: \texttt{Drowsy/Sleeping}, \texttt{Focused}, \texttt{Raising Hand}, \texttt{Using-Phone}. Tích hợp trực tiếp vào tập train Behavior Analysis.

    \item \textbf{MRL Eye Dataset} \href{https://www.kaggle.com/datasets/akashshingha850/mrl-eye-dataset}{[Kaggle]} --- \textbf{84.898 ảnh} trạng thái mắt (open/closed) từ 37 subjects. Dùng để train eye-state classifier phục vụ phát hiện \texttt{drowsy} qua Eye Aspect Ratio (EAR).
\end{itemize}

\vspace{0.1cm}
\begin{center}
\begin{tabularx}{0.96\textwidth}{lXr}
\toprule
\textbf{Nguồn} & \textbf{Mục đích sử dụng} & \textbf{Số mẫu} \\
\midrule
Self-collected   & Person detection \& behavior (góc thực tế FPT)           & 1.335 frames \\
SCB-Dataset      & Behavior classification, fine-tuning (20 classes)         & $\sim$7.428 ảnh \\
Ambient Classroom & Behavior detection, YOLOv8 format                        & $\sim$1.350 ảnh \\
MRL Eye          & Drowsiness detection qua eye state (EAR)                  & $\sim$84.898 ảnh \\
\bottomrule
\end{tabularx}
\end{center}

\end{document}
